\section{Cambiar de opción: borrar, dejar en pausa y retomar}
Solo puede haber una modalidad activa. Al final de tu trámite está el pliegue \emph{¿Elegiste otra modalidad o quieres dejar este trámite?}. Lo que pasa depende de si alguien más ya intervino:
\begin{itemize}
  \item \textbf{Nadie más intervino:} el trámite \textbf{se borra} definitivamente. La pantalla dice qué se borrará y pide confirmarlo.
  \item \textbf{Ya intervino alguien} (tu asesor firmó, el comité revisó, Titulación validó…): el trámite \textbf{queda en pausa} en \emph{Abandonadas}. Se conserva todo; mientras tanto nadie puede firmar ni revisar nada en él.
\end{itemize}
\captura{comun_pasante_01_dejar_pliegue}{Al final de tu trámite: «¿Elegiste otra modalidad o quieres dejar este trámite?».}
\captura{comun_pasante_02_borrar}{Si nadie más intervino, el trámite se borra: la pantalla dice qué se borrará y pide confirmarlo.}
\captura{comun_pasante_03_borrado}{Trámite borrado: vuelves a «¿Cómo quieres titularte?» para elegir otra opción.}
\subsection*{Iniciar otra opción teniendo una activa}
Si ya tienes una opción en proceso y entras a otra, la página te lo recuerda. Al pulsar iniciar, la plataforma explica qué pasará con la que tienes (se borra o queda en pausa) y te pide confirmarlo.
\captura{comun_pasante_04_otra_activa}{Si ya tienes una opción en proceso, la página de otra opción te lo recuerda antes de iniciar.}
\captura{comun_pasante_05_una_activa}{Solo puede haber una modalidad activa: antes de iniciar otra, la plataforma explica qué pasará con la que tienes y pide confirmarlo.}
\subsection*{Retomar un trámite en pausa}
En \emph{Mi trámite de licenciatura}, abre \emph{Abandonadas} y pulsa \boton{Retomar}. Vuelve justo donde se quedó; si tienes otra opción activa, antes se te explica qué pasará con ella.
\captura{comun_pasante_06_abandonadas}{«Abandonadas»: el trámite en pausa conserva todo; desde aquí lo retomas.}
\captura{comun_pasante_07_retomar}{Retomar: confirmas, y lo que tengas activo se borra o queda en pausa según el caso.}
\captura{comun_pasante_08_retomado}{El trámite retomado sigue justo donde se quedó.}
\captura{comun_pasante_09_pausa}{Si otras personas ya intervinieron, el trámite no se borra: queda en pausa y la pantalla explica por qué y qué se conserva.}
